% =====================================================================
\begin{thebibliography}{99}
\addcontentsline{toc}{chapter}{Bibliografía}
\markboth{Bibliografía}{}
\small

\bibitem{baurstrassen}
W.~Baur y V.~Strassen.
\newblock The complexity of partial derivatives.
\newblock \emph{Theoretical Computer Science}, 22(3):317--330, 1983.

\bibitem{baydin}
A.~G. Baydin, B.~A. Pearlmutter, A.~A. Radul y J.~M. Siskind.
\newblock Automatic differentiation in machine learning: a survey.
\newblock \emph{Journal of Machine Learning Research}, 18(153):1--43, 2018.
\newblock \url{https://jmlr.org/papers/v18/17-468.html}

\bibitem{chen}
T.~Chen, B.~Xu, C.~Zhang y C.~Guestrin.
\newblock Training deep nets with sublinear memory cost.
\newblock arXiv:1604.06174, 2016.
\newblock \url{https://arxiv.org/abs/1604.06174}

\bibitem{cmu11785}
B.~Raj y equipo docente de 11-785.
\newblock 11-785: Introduction to Deep Learning. Tarea~1, parte complementaria, «An Introduction to Autograd».
\newblock Carnegie Mellon University, otoño de 2021.
\newblock \url{https://github.com/CMU-IDeeL/new_grad}

\bibitem{cs231n}
A.~Karpathy y equipo docente de CS231n.
\newblock CS231n: Deep Learning for Computer Vision. Notas del curso, «Backpropagation, Intuitions».
\newblock Stanford University.
\newblock \url{https://cs231n.github.io/optimization-2/}

\bibitem{cs336a1}
Equipo docente de CS336.
\newblock CS336: Language Modeling from Scratch. Tarea~1, «Basics», versión 26.0.3.
\newblock Stanford University, primavera de 2026.
\newblock \url{https://github.com/stanford-cs336/assignment1-basics}

\bibitem{cs336a2}
Equipo docente de CS336.
\newblock CS336: Language Modeling from Scratch. Tarea~2, «Systems», versión 26.1.3.
\newblock Stanford University, primavera de 2026.
\newblock \url{https://github.com/stanford-cs336/assignment2-systems}

\bibitem{gemini}
Gemini Team, Google.
\newblock Gemini: a family of highly capable multimodal models.
\newblock arXiv:2312.11805, 2023.
\newblock \url{https://arxiv.org/abs/2312.11805}

\bibitem{gbc}
I.~Goodfellow, Y.~Bengio y A.~Courville.
\newblock \emph{Deep Learning}.
\newblock MIT Press, 2016.
\newblock §6.5, «Back-Propagation and Other Differentiation Algorithms».
\newblock \url{https://www.deeplearningbook.org/}

\bibitem{gw}
A.~Griewank y A.~Walther.
\newblock \emph{Evaluating Derivatives: Principles and Techniques of Algorithmic Differentiation}, 2.ª ed.
\newblock SIAM, 2008.

\bibitem{grok}
xAI.
\newblock Grok-1: código de ejemplo en JAX para cargar y ejecutar el modelo de pesos abiertos.
\newblock 2024.
\newblock \url{https://github.com/xai-org/grok-1}

\bibitem{grosse}
R.~Grosse.
\newblock Lecture~6: Automatic Differentiation. Notas del curso CSC421/2516, «Neural Networks and Deep Learning».
\newblock University of Toronto, 2019.
\newblock \url{https://www.cs.toronto.edu/~rgrosse/courses/csc421_2019/}

\bibitem{kaplan}
J.~Kaplan, S.~McCandlish, T.~Henighan, T.~B. Brown, B.~Chess, R.~Child, S.~Gray, A.~Radford, J.~Wu y D.~Amodei.
\newblock Scaling laws for neural language models.
\newblock arXiv:2001.08361, 2020.
\newblock \url{https://arxiv.org/abs/2001.08361}

\bibitem{karpathy}
A.~Karpathy.
\newblock micrograd: un pequeño motor de diferenciación automática en modo inverso sobre un grafo construido dinámicamente.
\newblock 2020.
\newblock \url{https://github.com/karpathy/micrograd}

\bibitem{linnainmaa}
S.~Linnainmaa.
\newblock The representation of the cumulative rounding error of an algorithm as a Taylor expansion of the local rounding errors.
\newblock Tesis de maestría, Universidad de Helsinki, 1970.

\bibitem{llama}
Meta.
\newblock Llama 2: código de inferencia de referencia (PyTorch).
\newblock 2023.
\newblock \url{https://github.com/meta-llama/llama}

\bibitem{mit18337}
C.~Rackauckas.
\newblock Parallel Computing and Scientific Machine Learning (18.337). Notas de la clase~8, «Forward-Mode Automatic Differentiation (AD) via High Dimensional Algebras».
\newblock Massachusetts Institute of Technology, 2020.
\newblock \url{https://book.sciml.ai/}

\bibitem{mit18s096}
A.~Edelman y S.~G. Johnson.
\newblock 18.S096: Matrix Calculus for Machine Learning and Beyond. Clase~5, partes~2 y~3.
\newblock Massachusetts Institute of Technology, MIT OpenCourseWare, IAP de enero de 2023.
\newblock \url{https://ocw.mit.edu/courses/18-s096-matrix-calculus-for-machine-learning-and-beyond-january-iap-2023/}

\bibitem{paszke}
A.~Paszke \emph{et al.}
\newblock PyTorch: an imperative style, high-performance deep learning library.
\newblock En \emph{Advances in Neural Information Processing Systems (NeurIPS)}, 2019.
\newblock \url{https://arxiv.org/abs/1912.01703}

\bibitem{pearlmutter}
B.~A. Pearlmutter.
\newblock Fast exact multiplication by the Hessian.
\newblock \emph{Neural Computation}, 6(1):147--160, 1994.

\bibitem{pytorch}
PyTorch.
\newblock Autograd mechanics. Documentación de PyTorch, versión 2.14.
\newblock \url{https://docs.pytorch.org/docs/stable/notes/autograd.html}

\bibitem{rumelhart}
D.~E. Rumelhart, G.~E. Hinton y R.~J. Williams.
\newblock Learning representations by back-propagating errors.
\newblock \emph{Nature}, 323:533--536, 1986.

\bibitem{d2l}
A.~Zhang, Z.~C. Lipton, M.~Li y A.~J. Smola.
\newblock \emph{Dive into Deep Learning}, versión 1.0.3.
\newblock Cambridge University Press, 2023.
\newblock §2.5, «Automatic Differentiation», y §5.3, «Forward Propagation, Backward Propagation, and Computational Graphs».
\newblock \url{https://d2l.ai/}

\end{thebibliography}
